\section{Task Description}
\label{sec:task_description}

India’s judicial system operates within the common law framework, where judges deliberate cases based on three fundamental pillars: (i)~the factual context of the case, (ii)~applicable statutory provisions, and (iii)~relevant judicial precedents. Our task is designed to simulate such realistic legal decision-making by leveraging Retrieval-Augmented Generation (RAG), enabling models to access external legal knowledge during inference.

% \noindent
Figure~\ref{fig:task-framework} illustrates our Legal Judgment Prediction (LJP) pipeline enhanced with RAG. The process begins with a full legal judgment document, which undergoes summarization to reduce its length and retain essential factual meaning. This is necessary because legal judgments tend to be long, and appending retrieved knowledge further increases the input size. Given limited model capacity and computational resources, we employ a summarization step (using \texttt{Mixtral-8x7B-Instruct-v0.1}) to create a condensed representation of both the input case and the retrieved legal context.

% \noindent
\paragraph{Prediction Task:} Based on the summarized factual description and the retrieved top-$k$ similar legal documents (statutes or precedents), the model predicts the likely court judgment. The prediction label \(y \in \{0, 1\}\) indicates whether the appeal is fully rejected ($0$) or fully/partially accepted ($1$). This binary framing captures the most common forms of judicial decisions in Indian appellate courts.

% \noindent
\paragraph{Explanation Task:} Alongside the decision, the model is also required to generate an explanation that justifies its output. This explanation should logically incorporate the facts, cited statutes, and relevant precedents retrieved during the RAG process. This step emulates how judges provide reasoned opinions in written judgments.

% \noindent
By structuring the LJP task in this manner, and by summarizing long documents and integrating retrieval-based augmentation, we study the effectiveness of RAG in producing judgments that are both faithful to legal reasoning and grounded in precedent and statute. The overall framework allows us to approximate real-world decision-making environments within Indian courtrooms.
